% ICML 2026 submission draft. Every table and every inline number is generated by make_tables.py
% from the committed result files (tables/numbers.tex, tables/*.tex); nothing is typed by hand.
\documentclass{article}
\usepackage{microtype,graphicx,booktabs,hyperref,amsmath,amssymb,amsfonts,amsthm}
\usepackage{algorithm}
\usepackage{algorithmic}
\usepackage{icml2026}
% generated by make_tables.py; do not edit
\newcommand{\NRegimes}{30}
\newcommand{\NSeedsPerRegime}{1{,}000}
\newcommand{\NRunsPerArm}{30{,}000}
\newcommand{\JNever}{71.53}
\newcommand{\JWald}{68.47}
\newcommand{\JFthreeTwenty}{69.43}
\newcommand{\SimVsNeverPct}{2.9}
\newcommand{\SimVsNeverDelta}{$-$2.10}
\newcommand{\SimVsNeverLo}{$-$2.22}
\newcommand{\SimVsNeverHi}{$-$1.98}
\newcommand{\SimVsWaldDelta}{+0.96}
\newcommand{\SimVsWaldLo}{+0.86}
\newcommand{\SimVsWaldHi}{+1.06}
\newcommand{\SimVsWaldPct}{1.4}
\newcommand{\SimVsWaldHiPct}{1.55}
\newcommand{\StressDelta}{$-$5.19}
\newcommand{\StressLo}{$-$5.40}
\newcommand{\StressHi}{$-$4.98}
\newcommand{\StressPct}{7.4}
\newcommand{\OrigDelta}{+2.50}
\newcommand{\OrigPct}{3.7}
\newcommand{\CTwoMaxHi}{0.006}
\newcommand{\COneMaxHi}{0.004}
\newcommand{\COneTotal}{0}
\newcommand{\CTwoTotal}{2}
\newcommand{\CThreeMinWorst}{64}
\newcommand{\CThreeMaxWorst}{99}
\newcommand{\WorstWald}{0.852}
\newcommand{\WorstLoWald}{0.829}
\newcommand{\WorstHiWald}{0.873}
\newcommand{\PooledWald}{6.6}
\newcommand{\WorstCusum}{0.987}
\newcommand{\WorstLoCusum}{0.978}
\newcommand{\WorstHiCusum}{0.992}
\newcommand{\PooledCusum}{7.9}
\newcommand{\WorstGate}{0.635}
\newcommand{\WorstLoGate}{0.605}
\newcommand{\WorstHiGate}{0.664}
\newcommand{\PooledGate}{4.3}
\newcommand{\WorstCusumRb}{0.830}
\newcommand{\WorstLoCusumRb}{0.805}
\newcommand{\WorstHiCusumRb}{0.852}
\newcommand{\PooledCusumRb}{4.4}
\newcommand{\WorstStop}{0.841}
\newcommand{\WorstLoStop}{0.817}
\newcommand{\WorstHiStop}{0.862}
\newcommand{\PooledStop}{4.2}
\newcommand{\CThreeMinLo}{0.605}
\newcommand{\AblNoPending}{0.940}
\newcommand{\AblLoNoPending}{0.924}
\newcommand{\AblHiNoPending}{0.953}
\newcommand{\AblNoRollback}{0.862}
\newcommand{\AblLoNoRollback}{0.839}
\newcommand{\AblHiNoRollback}{0.882}
\newcommand{\AblNoConf}{0.341}
\newcommand{\AblLoNoConf}{0.312}
\newcommand{\AblHiNoConf}{0.371}
\newcommand{\StressExceedWald}{23}
\newcommand{\OrigExceedWald}{2}
\newcommand{\StressExceedCusum}{17}
\newcommand{\OrigExceedCusum}{6}
\newcommand{\StressExceedGate}{13}
\newcommand{\OrigExceedGate}{2}
\newcommand{\StressExceedCusumRb}{16}
\newcommand{\OrigExceedCusumRb}{2}
\newcommand{\StressExceedStop}{17}
\newcommand{\OrigExceedStop}{1}
\newcommand{\RealRunsPerArm}{4{,}600}
\newcommand{\RealWaldExceedClean}{475}
\newcommand{\RealStopExceedClean}{311}
\newcommand{\RealWaldMaxClean}{175}
\newcommand{\RealShieldViolClean}{0}
\newcommand{\RealPiVsNeverCleanHfourteen}{75}
\newcommand{\RealPiVsNeverCleanHtwentyfour}{87}
\newcommand{\RealPiVsNeverCleanHfifty}{102}
\newcommand{\RealPiVsNeverCleanHhundred}{119}
\newcommand{\RealWaldExceedDropout}{890}
\newcommand{\RealStopExceedDropout}{886}
\newcommand{\RealWaldMaxDropout}{529}
\newcommand{\RealShieldViolDropout}{0}
\newcommand{\RealPiVsNeverDropoutHfourteen}{76}
\newcommand{\RealPiVsNeverDropoutHtwentyfour}{89}
\newcommand{\RealPiVsNeverDropoutHfifty}{105}
\newcommand{\RealPiVsNeverDropoutHhundred}{122}
\newcommand{\RealWaldExceedDelay}{1{,}017}
\newcommand{\RealStopExceedDelay}{979}
\newcommand{\RealWaldMaxDelay}{400}
\newcommand{\RealShieldViolDelay}{0}
\newcommand{\RealPiVsNeverDelayHfourteen}{74}
\newcommand{\RealPiVsNeverDelayHtwentyfour}{85}
\newcommand{\RealPiVsNeverDelayHfifty}{99}
\newcommand{\RealPiVsNeverDelayHhundred}{110}
\newcommand{\RealJNever}{597.89}
\newcommand{\RealJWald}{403.76}
\newcommand{\RealVsNeverMin}{75}
\newcommand{\RealVsNeverMax}{119}
\newcommand{\RealMaxExceed}{1{,}017}
\newcommand{\HarmDropMin}{33}
\newcommand{\HarmDropMax}{54}
\newcommand{\HarmDelayMin}{46}
\newcommand{\HarmDelayMax}{87}
\newcommand{\RegimeTestsSupported}{6}
\newcommand{\RegimeTestsTotal}{6}
\newcommand{\KeepLo}{88}
\newcommand{\KeepHi}{96}
\newcommand{\SpearLo}{0.69}
\newcommand{\SpearHi}{0.83}
\newcommand{\LargeXMin}{0.98}
\newcommand{\NUnits}{502}
\newcommand{\PThreeCells}{12}
\newcommand{\PThreeAttain}{12}
\newcommand{\PFourCells}{24}
\newcommand{\PFourAbove}{5}
\newcommand{\PFourMaxZ}{0.7}
\newcommand{\PFourTightEmp}{0.203}
\newcommand{\PFourTightBound}{0.198}
\newcommand{\PFourTightEmpB}{0.035}
\newcommand{\PFourTightBoundB}{0.034}
\newcommand{\VThreeNCand}{12}
\newcommand{\VThreeBest}{395.5}
\newcommand{\DefXTwelve}{0.25}
\newcommand{\DefXTwenty}{0.11}
\newcommand{\DefXThirty}{0.06}
\newcommand{\DefR}{0.3}
\newcommand{\DefD}{5}
\newcommand{\RealSixVsNever}{6.9}
\newcommand{\HarmNeverMin}{13}
\newcommand{\HarmNeverMax}{59}


\theoremstyle{plain}
\newtheorem{lemma}{Lemma}
\newtheorem{theorem}{Theorem}
\newtheorem{proposition}{Proposition}
\theoremstyle{definition}
\newtheorem{definition}{Definition}
\newcommand{\PP}{\mathbb{P}}
\newcommand{\E}{\mathbb{E}}
\newcommand{\one}{\mathbb{1}}

\icmltitlerunning{A Pathwise Harm Budget for Reversible Model Promotion}

\begin{document}
\twocolumn[
\icmltitle{A Pathwise Harm Budget for Reversible Model Promotion\\
under Delayed, Missing and Scarce Feedback}
\icmlsetsymbol{equal}{*}
\begin{icmlauthorlist}
\icmlauthor{Anonymous Authors}{anon}
\end{icmlauthorlist}
\icmlaffiliation{anon}{Anonymous Institution}
\icmlcorrespondingauthor{Anonymous Authors}{anon@example.com}
\icmlkeywords{safe deployment, model promotion, confidence sequences, delayed feedback, conservative bandits}
\vskip 0.3in
]
\printAffiliationsAndNotice{}

\begin{abstract}
Replacing a deployed model with a challenger is an intervention whose harm is learned only through labels that arrive late, arrive for a sampled subset of rounds, or never arrive.
Certifying that a switch will repay its cost requires a drift assumption.
We bound instead the harm the intervention actually causes.
For a user-chosen budget $R$ and level $\delta$, we require that with probability at least $1-\delta$ the cumulative excess loss against never switching, plus switching fees, stays below $R$ at every round, for every loss sequence and every delay or missing-label mechanism.
A shield around any promotion rule achieves this from a design-based inverse-propensity bound on the realised harm and two reserves, one for the fee of a rollback and one for feedback still in transit.
A matching lower bound shows that both reserves are necessary, and in the feedback-blind window the shield holds a first promotion exactly as long as any valid controller can.
In a registered simulation of \NRegimes{} regimes (\NRunsPerArm{} runs per arm) the shield has \COneTotal{} violations for full sampling and at most one per \NRunsPerArm{} runs with paced sampling, while tuned Wald, CUSUM, e-process-gate, CUSUM-rollback and stop-loss rules exceed the same budget in \CThreeMinWorst--\CThreeMaxWorst\% of runs in their worst regime.
On 14 real streams under three feedback conditions the shield has \RealShieldViolClean{} violations in \RealRunsPerArm{} runs per arm and condition, where tuned Wald+cutoff exceeds the same budget in up to \RealMaxExceed{} runs.
Among controllers that keep the budget, the shield improves on never switching by \SimVsNeverPct\% pooled in simulation and by \RealVsNeverMin--\RealVsNeverMax{} cost units per episode on real streams, and it retains \KeepLo--\KeepHi\% of tuned Wald+cutoff's gain whenever one delay window holds less than half the budget headroom, as the bound predicts.
\end{abstract}

\section{Introduction}
A team that has trained a challenger model must decide when to let it replace the model in production.
The decision is made with feedback that is slow, partial and sometimes absent: labels arrive after a delay, only a purchased subset of rounds is labelled, and some labels never arrive.
Sequential tests and change detectors answer a different question from the one the team faces.
An alarm or a certificate that the challenger is better now says nothing about what the switch will have cost by the time the labels come in, and a certificate that the switch will pay for itself over a future horizon cannot be had without a drift assumption: two worlds that agree on every visible label, one in which the challenger's advantage persists and one in which it reverses, are indistinguishable, so no rule separates them.

We take the other route and bound what the intervention does rather than what it will do.
Realised harm, the excess loss of the deployed action against never switching plus the switching fees paid, is observable in hindsight from sampled labels, and the intervention is reversible, so a controller can stop the harm once it sees it.
The target is a risk budget: for a chosen $R$ and $\delta$, with probability at least $1-\delta$ the harm never exceeds $R$ at any round (Section~\ref{sec:setting}).
The statement is pathwise and design-based.
It holds for every loss sequence, including arbitrary drift, reversals and dependence, and for any delay or missing-label mechanism, because all probability comes from the controller's own query coins.

The difficulty is that harm becomes known only after its delay.
A controller that rolls back once observed harm reaches the budget has already incurred the harm of the rounds whose labels are in transit, and the fee of the rollback itself.
We show that a shield around any promotion rule meets the budget by holding two reserves: the rollback fee and the worst case of feedback still pending (Section~\ref{sec:method}).
A matching lower bound shows that both reserves are necessary and that the shield is exactly tight on the all-harm tape (Section~\ref{sec:theory}).
The same accounting predicts how much of an unconstrained rule's gain a valid controller can retain, in terms of the exposure per delay window relative to the headroom $h=R-2k$.

\paragraph{Contributions.}
(1) A formulation: a pathwise, high-probability budget on realised intervention harm under delayed, missing and sampled feedback, with switching fees and reversibility.
(2) A method: a shield around any promotion rule, built on an anytime inverse-propensity bound (Bennett and exact coin mixtures, with settlement of pending rounds), with Theorem~\ref{thm:main}.
(3) Theory: a lower bound showing the two reserves are necessary and attained (Proposition~\ref{prop:reserves}), and a ruin bound for the first deployment spell that is tight at unit delay (Proposition~\ref{prop:price}).
(4) A predicted regime map in $x=rd/h$, confirmed on all \RegimeTestsTotal{} pre-stated tests.
(5) Registered confirmation of the budget on \NRegimes{} simulated regimes, and a frozen-configuration study on 14 real streams under three feedback conditions, against five heuristic or unconstrained comparators, all of which break the budget.

\section{Setting and target}\label{sec:setting}
Rounds are $t=1,\dots,T$. Two frozen models produce losses $\ell_t(0),\ell_t(1)\in[0,1]$, with action $0$ the incumbent; $Z_t=\ell_t(0)-\ell_t(1)$. A flag $g_t\in\{0,1\}$, observed before acting, marks rounds on which the models disagree, and $Z_t=0$ when $g_t=0$. The controller picks $a_t\in\{0,1\}$ with $a_0=0$, pays a fee $k$ per action change ($N_t$ changes through $t$), and may buy a label for round $s$ at price $c_q$ out of a budget $B$. A label reveals both losses of round $s$. It is released at the start of round $s+d_s$ for some $d_s\ge1$, or never. Let
\begin{equation*}
\chi_t=-a_tg_tZ_t,\qquad H_t=kN_t+\sum_{s\le t}\chi_s .
\end{equation*}
$\chi_t$ is the excess loss of the deployed action over action $0$, and $H_t$ is the intervention harm. Gains lower $H_t$, so a challenger that earns its keep frees budget. Label prices are accounted separately through $B$. A round with $a_t=1,g_t=1$ is \emph{exposed}.

\paragraph{Observation model.}
(O1) Outcomes are exogenous: actions and queries do not change losses. (O2) On an exposed round the controller chooses $\pi_t\in\{0\}\cup[\pi_{\min},1]$ from its observed history and queries with a private coin $q_t\sim\mathrm{Bernoulli}(\pi_t)$, independent of the tape given the past; $\pi_t=0$ is a \emph{blind} round. (O3) Delays are arbitrary: random, outcome-dependent, adversarial, or infinite, provided they are fixed by the tape or by the past and not by the current coin. (O4) Losses lie in $[0,1]$. Nothing is assumed about the law of $(Z_t)$.

\begin{definition}[$(R,\delta)$ risk budget]
A controller satisfies the budget if $\PP(\sup_{t\le T}H_t>R)\le\delta$ for every tape satisfying O1--O4, where the probability is over the query coins.
\end{definition}
With $\pi_t=1$ on every monitored round the budget holds with $\delta=0$.

\section{The shield}\label{sec:method}
\paragraph{Anytime bound on harm.}
Split exposed rounds into blind rounds (count $n_{\rm b}$) and monitored rounds ($\pi_s>0$). After releases at the start of round $t+1$ define, over rounds $s\le t$,
{\small\begin{align*}
A&=\sum_{\substack{s\ \text{queried,}\\ \text{released}}}\!\chi_s/\pi_s,\quad
P=\sum_{\substack{s\ \text{queried,}\\ \text{pending}}}\!1/\pi_s,\\
V&=\sum_{s\ \text{monitored}}(1/\pi_s-1),\quad
M_t=\sum_{s\ \text{monitored}}\chi_s(1-q_s/\pi_s).
\end{align*}}
Since $\chi_s\le1$, $\sum_{s\le t}\chi_s\le n_{\rm b}+A+P+M_t$. The term $M_t$ is a martingale in the query coins.

\begin{lemma}[Anytime width]\label{lem:bennett}
Let $b=\max(1,1/\pi_{\min}-1)$, $\phi(\lambda)=(e^{\lambda b}-1-\lambda b)/b^2$, and $\Lambda$ a finite grid with weights $w_\lambda$ summing to one. With probability at least $1-\delta$, for all $t$ simultaneously, $M_t\le D_\delta(V_t)$, where $D_\delta(0)=0$ and $D_\delta(V)=\min_{\lambda\in\Lambda}[\ln(1/(\delta w_\lambda))+\phi(\lambda)V]/\lambda$.
\end{lemma}
Hence $U_t=kN_t+n_{\rm b}+A+P+D_\delta(V_t)$ bounds $H_t$ at every $t$ with probability $1-\delta$. The proof is Bennett's inequality for the increments $\chi_s(1-q_s/\pi_s)$, a mixture over $\lambda$ and Ville's inequality (Appendix~\ref{app:proofs}). Appendix~\ref{app:v2} gives a sharper version (Lemma~\ref{lem:coin}): since $\E\exp(\lambda x(1-q/\pi))$ is convex in $x\in[-1,1]$, it is at most $e^{\psi(\lambda,\pi)}$ with $\psi$ the larger of its two endpoint values, and mixing $e^{\lambda M_t-\sum_s\psi(\lambda,\pi_s)}$ over a small grid gives $M_t\le D_t$ for any per-round $\pi_s\in[\pi_{\min},1]$, with no global $b$. The paced variant uses it. Confidence sequences of this type are standard \cite{howard2021time,waudbysmith2024betting}; the specifics here are the target, the design-based validity under dependence and arbitrary delays, and the reserves below.

\paragraph{Shield.}
The shield wraps an arbitrary inner promotion rule (here tuned Wald+cutoff). When the inner rule proposes $a_t=1$, the shield evaluates, for a monitoring rate $\pi$ (cost $\kappa(\pi)=1/\pi$, or $1$ if blind),
{\small\begin{align*}
W_t(\pi)&=k\,(N_{t-1}+\one\{a_{t-1}=0\})+\underbrace{k}_{\text{rollback}}+n_{\rm b}+A+P\\
&\quad+\underbrace{g_t\kappa(\pi)}_{\text{this round}}+D_\delta\big(V+g_t(1/\pi-1)\big),
\end{align*}}
with $A,P,V,n_{\rm b}$ evaluated at the start of round $t$. It deploys $a_t=1$ iff $W_t(\pi_t)\le R$ for the cheapest $\pi_t$ in a menu such as $\{\pi_{\rm lo},1\}$, falling back to $\pi=1$ and then to $a_t=0$, which is a rollback if $a_{t-1}=1$. The two reserves are the rollback fee $k$, and the worst case $\kappa(\pi)$ for the current round together with $P$ for the rounds whose labels have not arrived. If the remaining horizon cannot exhaust the budget ($kN_{t-1}+n_{\rm b}+A+P+D_\delta(V)+T-t+1\le R$ while $a_{t-1}=1$) the shield rides out to $T$ without queries and without the rollback reserve. A fraction of the label budget is held back while $a=0$ so that monitoring labels exist after a promotion. The paced variant sets $\pi$ from the labels remaining per expected exposed round and books sampled rounds at worst case until a settlement window has passed (Appendix~\ref{app:v2}). Algorithm~\ref{alg:shield} summarises one round.

\begin{algorithm}[t]
\caption{One round of the risk-budget shield}\label{alg:shield}
\begin{algorithmic}
\STATE \textbf{Input:} inner proposal $\tilde a_t$, budget $R$, level $\delta$, menu of rates
\STATE Release labels; update $A,P,V,n_{\rm b},N_{t-1}$
\IF{$\tilde a_t=0$} \STATE $a_t\leftarrow0$ \COMMENT{rollback fee already reserved}
\ELSE
  \STATE choose the cheapest $\pi$ in the menu with $W_t(\pi)\le R$ (else $\pi=1$)
  \IF{$W_t(\pi)\le R$} \STATE $a_t\leftarrow1$; if $g_t=1$ draw $q_t\sim\mathrm{Bern}(\pi)$ (blind if $\pi=0$)
  \ELSE \STATE $a_t\leftarrow0$ \ENDIF
\ENDIF
\end{algorithmic}
\end{algorithm}

\section{Guarantees}\label{sec:theory}
\begin{theorem}[Pathwise risk budget]\label{thm:main}
Under O1--O4, for any inner policy, any tape and any $R\ge0$, the shield satisfies $\PP(\sup_tH_t>R)\le\delta$. If $\pi_t=1$ on every monitored round, $\sup_tH_t\le R$ surely.
\end{theorem}
\begin{proof}
Let $E$ be the event of Lemma~\ref{lem:bennett}, which has probability at least $1-\delta$ (it is sure when $\pi\equiv1$). Fix a round $t$ with $a_t=1$. The fees paid through $t$ are $k(N_{t-1}+\one\{a_{t-1}=0\})$. The excess loss of rounds $s\le t-1$ is at most $n_{\rm b}+A+P+M_{t-1}$, round $t$ adds at most $g_t\kappa(\pi_t)$ to $n_{\rm b}+A+P$ and $g_t(1/\pi_t-1)$ to $V$, and $M\le D_\delta(V)$ on $E$ with $D_\delta$ nondecreasing. Hence $H_t+k\le W_t(\pi_t)\le R$ on $E$, so $H_t\le R-k$ at every deployed round. For an arbitrary $t$ let $t'\le t$ be the last deployed round, if any; otherwise $H_t=0$. Between $t'$ and $t$ the only change in $H$ is the rollback fee $k$ at $t'+1$, since $\chi_s=0$ when $a_s=0$ and there is no later switch. So $H_t\le R-k+k=R$. The argument needs only that some valid upper bound on $H_t+k$ is verified at each deployed round, so it also covers the sharper bound of Appendix~\ref{app:v2}, and the ride-out branch (Appendix~\ref{app:proofs}).
\end{proof}

\begin{proposition}[Both reserves are necessary]\label{prop:reserves}
Suppose every label requested at round $s$ is released no earlier than $s+d$, $d\ge1$, and a controller satisfies the $(R,\delta)$ budget on every tape. Fix an environment $\mathcal E$ and rounds $j$, $m$ with $1\le m\le d$. Let $\mathcal A$ be the event that the first promotion is at round $j$ and $a_s=1$ for $s=j,\dots,j+m-1$; let $e$ and $n$ be the numbers of disagreement rounds in $[j,j+m-1]$ and $[j,T]$. If $\PP_{\mathcal E}(\mathcal A)>\delta$ then $R\ge\min(2k+e,\,k+n)$.
\end{proposition}
\begin{proof}
Let $\mathcal E'$ equal $\mathcal E$ before round $j$ and have $Z_s=-1$ on every later disagreement round, with the same flags. A label visible at a round $s\le j+d-1$ was requested at or before $s-d\le j-1$, where the two environments coincide, so with the same coins the actions agree through round $j+d-1$. Since $m\le d$, $\mathcal A$ depends only on those actions and $\PP_{\mathcal E'}(\mathcal A)=\PP_{\mathcal E}(\mathcal A)>\delta$. Under $\mathcal E'$ every $\chi_s\ge0$, so $H_t$ never decreases. On $\mathcal A$, $H_{j+m-1}=k+e$. If the controller later rolls back at $\rho$, $H_\rho\ge2k+e$; if it never does, $H_T=k+n$. As $n\ge e$, $\sup_tH_t\ge\min(2k+e,k+n)$ on $\mathcal A$, and validity forces $R\ge\min(2k+e,k+n)$.
\end{proof}

The three terms of $W_t$ on a first promotion are exactly these: $k$ for the promotion, $k$ reserved for the rollback, and one unit per exposed round whose label cannot yet have arrived. On an all-disagreement, all-harm tape the full-sampling shield holds a first promotion for exactly $\lfloor R-2k\rfloor$ exposed rounds, which is the most Proposition~\ref{prop:reserves} allows when $d\ge\lfloor R-2k\rfloor$. This was checked in \PThreeAttain{} of \PThreeCells{} $(d,R)$ cells (Appendix Table~\ref{tab:prop3}). The ride-out branch matches the $k+n$ term, using the remaining round count as an upper bound on $n$. The proposition concerns the feedback-blind window only.

\begin{proposition}[Ruin bound for a first spell]\label{prop:price}
Take the full-sampling shield with fixed delay $d$ during its first deployment spell, while labels remain and the inner rule keeps proposing the challenger. Let $y_1,y_2,\dots$ be the values of $Z$ on the exposed rounds in order, with $\E[y_i\mid\text{past}]\ge\mu>0$, and let $m=R-2k-d$. Then $\PP(\text{forced rollback in the spell})\le e^{-2\mu m}$. If $y_i\in\{-1,+1\}$ with $\PP(y_i=+1\mid\text{past})\ge p>\tfrac12$, the probability is at most $((1-p)/p)^{\lfloor m\rfloor+1}$.
\end{proposition}
\emph{Proof sketch.} During the spell $W_t\le2k+d-S_n$ with $S_n=y_1+\dots+y_n$, because the pending rounds and the current round number at most $d$ and the released rounds form a prefix of the exposed sequence. A forced rollback needs $S_n<-m$ for some $n$, and $e^{-2\mu S_n}$ (respectively $(q/p)^{S_n}$) is a nonnegative supermartingale, so Ville's inequality gives the bounds. \hfill$\square$

Proposition~\ref{prop:price} was checked by simulation in \PFourCells{} cells ($p\in\{.6,.7,.8\}$, $d\in\{1,5\}$, $R\in\{10,12,15,20\}$): all lie within Monte Carlo error of the bound, and at $d=1$ it is tight (\PFourTightEmp{} against \PFourTightBound{} at $p=.6,R=10$; \PFourTightEmpB{} against \PFourTightBoundB{} at $p=.7,R=10$). At $d=5$ it is looser because the pending count is usually below $d$ (Appendix Table~\ref{tab:prop4}).

\paragraph{What the bound predicts.}
Let $r$ be the disagreement rate while deployed and $h=R-2k$ the headroom. Proposition~\ref{prop:reserves} limits any valid controller to holding a first promotion for at most $h$ exposed rounds before any label has arrived, and one delay window contains about $rd$ of them. When $x=rd/h$ is small, a promotion can be held until labels arrive and the shield can behave like the unconstrained rule. When $x>1$, no valid controller can remain deployed through one delay window on the strength of the prior alone, and it must cycle (deploy, roll back, wait, re-enter) until gains banked from released labels cover the window. We therefore predict that the share $L$ of the unconstrained rule's improvement that the shield gives up increases with $x$, is small for $x<.5$, and is near one for $x>1$. The prediction was committed with its test script before the test was first run, and is evaluated in Section~\ref{sec:exp}.

\section{Experiments}\label{sec:exp}
\paragraph{Simulation.}
The registered confirmation (frozen before any confirmation seed was opened) uses the 24 regimes of the original family and 6 stress regimes built to stress delayed and missing feedback (adverse informative delay, dropout of unfavourable labels, a bait-and-reverse path, dependence $\rho=.95$), with $T=300$, label budget $B=60$, $k=3$, $c_q=.05$ and $\NSeedsPerRegime$ seeds per regime ($\NRunsPerArm$ runs per arm), run once. The inner rule is Wald+cutoff tuned by a 32-configuration random search on separate seeds, with the same search budget as every comparator. Budgets are $R\in\{12,20,30\}$, that is $h\in\{6,14,24\}$. The cost $J$ is deployed loss plus label, switching and compute fees. Comparators are tuned Wald+cutoff, two-sided CUSUM, a betting e-process gate over Wald, Wald+cutoff with a one-sided CUSUM rollback, and a naive stop-loss (Wald+cutoff that rolls back when fees plus resolved harm exceed $R-k$, with no pending reserve and no width). None is a native-source reproduction: they are in-house heuristics or documented adapters.

\paragraph{Real streams.}
We use 14 cached streams of paired 0/1 losses (230 episodes, $T=2000$, $B=200$, $k=5$, 20 policy seeds per episode), under the stream's own delay (clean), adverse dropout (unfavourable labels never released with probability $.6$) and adverse delay (unfavourable labels delayed by 400 further events); the two corrupted conditions change only the feedback. Budgets are $R=10+h$ with $h\in\{6,14,24,50,100\}$. Every arm is frozen from simulation and none was run or tuned on these streams before this study. The challenger is worse on two streams (airlines, bank) and better on twelve.

\begin{table*}[t]
\caption{Registered confirmation C1--C6 (simulation, \NRunsPerArm{} runs per arm). Source: \texttt{confirm\_summary.json}.}\label{tab:confirm}
\centering\small
\begin{tabular}{@{}l p{3.9cm} p{5.7cm} l@{}}
\toprule
ID & Registered statistic & Result & Verdict \\
\midrule
C1 & Full-sampling shield, $R\in\{12,20,30\}$, and blind-round variant: violations of $R$ & 0 / 30{,}000 in each of 4 arms; per-regime Wilson upper $\le 0.004$ & supported \\
C2 & Paced shield, $\delta=.05$, $R\in\{20,30\}$: violations & 1 and 1 / 30{,}000; largest per-regime Wilson upper 0.006 & supported \\
C3 & Five comparators, worst-regime exceedance of $R=20$ (Wilson lower bound) & all five above $.05$; smallest lower bound 0.605 & supported \\
C4 & Pooled $\Delta J$ of the $R=20$ shield, 95\% CI & vs never $-$2.10 [$-$2.22, $-$1.98]; vs Wald+cutoff +0.96 [+0.86, +1.06] & supported \\
C5 & $\Delta J$ vs Wald+cutoff, 6 stress regimes (24 original regimes) & $-$5.19 [$-$5.40, $-$4.98] (+2.50) & supported \\
C6 & Ablations at $R=12$, worst-regime violation rate (no pending / no rollback reserve / no width) & 0.940 / 0.862 / 0.341 & supported \\
\bottomrule
\end{tabular}

\end{table*}

\begin{table*}[t]
\caption{Simulation, $R=20$, pooled over \NRegimes{} regimes. $J$ is mean cost; ``vs never'' is the relative difference in $J$; harm quantiles are of $\sup_tH_t$. The last two columns give the exceedance rate of $R=20$ pooled and in the worst regime. The two shields keep the budget; the comparators do not.}\label{tab:sim}
\centering\scriptsize
\begin{tabular}{@{}l r r r r r r l l@{}}
\toprule
Controller & $J$ & vs never (\%) & labels & $q_{99}$ harm & max harm & pooled viol.\ (\%) & worst-regime rate [Wilson] & worst regime \\
\midrule
Never switch & 71.53 & +0.0 & 0.0 & 0 & 0 & 0.0 & 0.000 & -- \\
Wald+cutoff (tuned) & 68.47 & $-$4.3 & 38.8 & 37 & 59 & 6.6 & 0.852 [0.829, 0.873] & dropout\_adverse \\
CUSUM (tuned) & 68.97 & $-$3.6 & 47.0 & 34 & 117 & 7.9 & 0.987 [0.978, 0.992] & expensive\_switch \\
E-process gate (tuned) & 71.89 & +0.5 & 49.7 & 33 & 57 & 4.3 & 0.635 [0.605, 0.664] & dropout\_adverse \\
CUSUM rollback, $h=8$ & 68.62 & $-$4.1 & 49.8 & 35 & 71 & 4.4 & 0.830 [0.805, 0.852] & dropout\_adverse \\
Stop-loss, $R=20$ & 68.79 & $-$3.8 & 49.8 & 34 & 59 & 4.2 & 0.841 [0.817, 0.862] & dropout\_adverse \\
\midrule
Shield, full sampling, $R=20$ & 69.43 & $-$2.9 & 44.2 & 20 & 20 & 0.000 & 0.000 [0.000, 0.004] & -- \\
Shield, paced, $R=20$, $\delta=.05$ & 69.80 & $-$2.4 & 37.8 & 16 & 22 & 0.003 & 0.001 [0.000, 0.006] & bait\_reversal \\
\bottomrule
\end{tabular}

\end{table*}

\subsection{The budget holds and the comparators break it}
Table~\ref{tab:confirm} lists the six registered outcomes, all supported. The full-sampling shield and the blind-round variant had \COneTotal{} violations in each of four arms ($R\in\{12,20,30\}$ and the variant), the paced shield with $\delta=.05$ had \CTwoTotal{} violations in two arms, against the guaranteed rate of $.05$, and the largest per-regime Wilson upper bound over all six arms is \CTwoMaxHi. Table~\ref{tab:sim} shows the comparators at $R=20$: each exceeds the budget in the majority of runs in its worst regime (Wald+cutoff \WorstWald, CUSUM \WorstCusum, gate \WorstGate, CUSUM rollback \WorstCusumRb, stop-loss \WorstStop). The exceedance is concentrated in the stress regimes: pooled over the six stress regimes, Wald+cutoff exceeds $R=20$ in \StressExceedWald\% of runs and over the 24 original regimes in \OrigExceedWald\% (CUSUM \StressExceedCusum\% and \OrigExceedCusum\%; stop-loss \StressExceedStop\% and \OrigExceedStop\%; Appendix Table~\ref{tab:split}). The two comparators that try to keep a budget, the stop-loss and the CUSUM rollback, fail in the same regimes as the rules they wrap, because the harm of rounds whose labels are in transit and the rollback fee are invisible to them. The ablations confirm that each part of the accounting is needed: without the pending reserve the worst-regime violation rate at $R=12$ is \AblNoPending, without the rollback reserve \AblNoRollback, and without the confidence width \AblNoConf (Table~\ref{tab:ablation}).

Among controllers that keep the budget, the shield beats never switching: at $R=20$ the pooled cost is \JFthreeTwenty{} against \JNever{} for never switching, a difference of \SimVsNeverDelta{} with 95\% interval [\SimVsNeverLo, \SimVsNeverHi], or \SimVsNeverPct\%. In the six stress regimes, where delayed and missing feedback mislead the heuristics, the shield also has lower cost than tuned Wald+cutoff by \StressDelta{} ([\StressLo, \StressHi], \StressPct\%).

\begin{table*}[t]
\caption{Real streams, pooled over 14 streams (equal weight), \RealRunsPerArm{} runs per arm and condition. $J$ is mean cost per episode, ``vs never'' the difference in $J$, ``viol.'' the number of runs above the budget: $R=24$ for the comparators, the arm's own $R$ ($24$, $60$ or $110$) for the shield and stop-loss. Source: \texttt{real\_tables.md}.}\label{tab:real}
\centering\scriptsize
\begin{tabular}{@{}l r r r r r r r r r @{}}
\toprule
 & \multicolumn{3}{c}{Clean} & \multicolumn{3}{c}{Adverse dropout} & \multicolumn{3}{c}{Adverse delay} \\
\cmidrule(lr){2-4}\cmidrule(lr){5-7}\cmidrule(lr){8-10}
Controller & $J$ & vs never & viol. & $J$ & vs never & viol. & $J$ & vs never & viol. \\
\midrule
Never switch & 597.89 & 0.0 & 0 & 597.89 & 0.0 & 0 & 597.89 & 0.0 & 0 \\
Wald+cutoff & 403.76 & $-$194.1 & 475 & 406.62 & $-$191.3 & 890 & 418.71 & $-$179.2 & 1{,}017 \\
CUSUM & 418.42 & $-$179.5 & 709 & 414.03 & $-$183.9 & 889 & 426.68 & $-$171.2 & 1{,}050 \\
E-process gate & 419.26 & $-$178.6 & 504 & 407.44 & $-$190.4 & 644 & 421.29 & $-$176.6 & 1{,}142 \\
CUSUM rollback $h=8$ & 415.60 & $-$182.3 & 517 & 415.44 & $-$182.4 & 884 & 460.46 & $-$137.4 & 980 \\
Stop-loss, $R=24$ & 419.76 & $-$178.1 & 311 & 414.87 & $-$183.0 & 886 & 447.45 & $-$150.4 & 979 \\
\midrule
Shield, full, $R=24$ & 522.82 & $-$75.1 & 0 & 521.74 & $-$76.2 & 0 & 524.20 & $-$73.7 & 0 \\
Shield, paced, $R=24$ & 519.09 & $-$78.8 & 0 & 518.95 & $-$79.0 & 0 & 520.44 & $-$77.5 & 0 \\
Shield, full, $R=60$ & 495.84 & $-$102.0 & 0 & 492.82 & $-$105.1 & 0 & 498.64 & $-$99.2 & 0 \\
Shield, paced, $R=60$ & 450.48 & $-$147.4 & 0 & 447.24 & $-$150.7 & 0 & 451.15 & $-$146.7 & 0 \\
Shield, paced, $R=110$ & 433.85 & $-$164.0 & 0 & 431.58 & $-$166.3 & 0 & 438.72 & $-$159.2 & 0 \\
\bottomrule
\end{tabular}

\end{table*}

\begin{table}[t]
\caption{Exceedance rate of $R=20$ pooled (mean of per-regime rates) over the six stress regimes and the 24 original regimes.}\label{tab:split}
\centering\scriptsize
\begin{tabular}{@{}l r r@{}}
\toprule
Arm & 6 stress regimes & 24 original regimes \\
\midrule
wald\_cut & 0.231 & 0.025 \\
cusum & 0.168 & 0.056 \\
gate\_wald & 0.132 & 0.020 \\
cusum\_rb\_h8 & 0.160 & 0.015 \\
stoploss\_R20 & 0.168 & 0.011 \\
f3\_pi1\_R20 & 0.000 & 0.000 \\
f3v2d\_R20 & 0.000 & 0.000 \\
\bottomrule
\end{tabular}

\end{table}

\begin{table}[t]
\caption{Ablations at $R=12$ (registered C6): worst-regime violation rate with the Wilson 95\% interval.}\label{tab:ablation}
\centering\scriptsize
\begin{tabular}{@{}l l r l@{}}
\toprule
Ablation ($R=12$) & worst regime & rate & Wilson 95\% \\
\midrule
no pending reserve & dropout\_adverse & 0.940 & [0.924, 0.953] \\
no rollback reserve & dropout\_adverse & 0.862 & [0.839, 0.882] \\
no confidence width & bait\_reversal & 0.341 & [0.312, 0.371] \\
\bottomrule
\end{tabular}

\end{table}

\subsection{Real streams}
Table~\ref{tab:real} gives the pooled results. The shield had \RealShieldViolClean{} violations of its own budget in every arm, headroom and condition (\RealRunsPerArm{} runs each). Tuned Wald+cutoff exceeded $R=24$ in \RealWaldExceedClean{} runs on clean feedback, \RealWaldExceedDropout{} under adverse dropout and \RealWaldExceedDelay{} under adverse delay, with a largest harm of \RealWaldMaxClean, \RealWaldMaxDropout{} and \RealWaldMaxDelay{} against a budget of 24; the stop-loss broke its own $R=24$ in \RealStopExceedClean, \RealStopExceedDropout{} and \RealStopExceedDelay{} runs. The shield improves on never switching by \RealPiVsNeverCleanHfourteen{} to \RealPiVsNeverCleanHhundred{} cost units per episode with full sampling at $h=14$ to $100$ on clean feedback, and by more with paced sampling at large headroom. At $h=6$ the budget leaves almost no room (a promotion holds for $\lfloor R-2k\rfloor=6$ exposed rounds), and the shield is within \RealSixVsNever{} units of never switching.

Table~\ref{tab:harm} isolates the two streams on which the challenger is worse. Under corrupted feedback, the shield has lower cost than tuned Wald+cutoff by \HarmDropMin{} to \HarmDropMax{} units under dropout and \HarmDelayMin{} to \HarmDelayMax{} under delay, since the corrupted feedback withholds the labels that would reveal the harm and Wald+cutoff, which sees only released labels, keeps the challenger deployed. The shield remains above never switching there, by \HarmNeverMin{} to \HarmNeverMax{} units under corrupted feedback, since the budget permits some harm before the shield withdraws.

\begin{table}[t]
\caption{Two harmful streams (airlines, bank): $\Delta J$ of the full-sampling shield against Wald+cutoff and against never switching, by headroom $h$.}\label{tab:harm}
\centering\scriptsize
\resizebox{\columnwidth}{!}{\begin{tabular}{@{}l l r r r r r @{}}
\toprule
 & & \multicolumn{5}{c}{headroom $h$ ($R=10+h$)} \\
\cmidrule(lr){3-7}
Feedback & vs & 6 & 14 & 24 & 50 & 100 \\
\midrule
Clean & wald\_cut & +2.2 & +2.3 & +0.9 & +2.6 & +6.2 \\
 & never & +10.9 & +11.0 & +9.6 & +11.3 & +14.9 \\
Adverse dropout & wald\_cut & $-$54.5 & $-$52.6 & $-$51.5 & $-$42.0 & $-$32.9 \\
 & never & +13.4 & +15.3 & +16.4 & +25.9 & +35.0 \\
Adverse delay & wald\_cut & $-$87.5 & $-$84.3 & $-$81.2 & $-$66.8 & $-$45.5 \\
 & never & +16.9 & +20.0 & +23.2 & +37.6 & +58.8 \\
\bottomrule
\end{tabular}
}
\end{table}

\subsection{Regime map}
We tested the prediction of Section~\ref{sec:theory} on the label-scarce simulation family and on the ten real streams where Wald+cutoff gains at least 10\% of $J$(never) (\NUnits{} units of regime or stream, arm and headroom). Table~\ref{tab:regime} reports Spearman correlations between $x=rd/h$ and $L=(J_{\rm shield}-J_{\rm Wald})/(J_{\rm never}-J_{\rm Wald})$ with regime- or stream-bootstrap intervals, and the median of $L$ by bin of $x$. All \RegimeTestsSupported{} of \RegimeTestsTotal{} pre-stated tests are supported, with correlations from \SpearLo{} to \SpearHi. For the paced arms the median $L$ is small below $x=.5$, so the shield retains \KeepLo--\KeepHi\% of the unconstrained rule's improvement, and it is at least \LargeXMin{} for $x\ge1$, where the shield is about as costly as never switching, consistent with Proposition~\ref{prop:reserves}: once one window holds more exposure than the headroom, a valid controller cannot stay deployed through it without banked gains. Full sampling gives up more at small $x$ because it spends one label per exposed round. A regression of $L$ on $\log x$ and the log labels per exposed round gives a positive coefficient on $\log x$ in all six fits (Appendix Table~\ref{tab:ols}), so the effect is not explained by label use. In the registered simulation the default regimes have $r=\DefR$ and $d=\DefD$, so $x=\DefXTwelve$, \DefXTwenty{} and \DefXThirty{} at $R=12,20,30$, all in the retaining range. The pooled difference to Wald+cutoff at $R=20$ is \SimVsWaldDelta{} ([\SimVsWaldLo, \SimVsWaldHi], \SimVsWaldPct\% of its cost), and \OrigDelta{} on the 24 original regimes (\OrigPct\%). A label-scarce variant (v3) selected on simulation narrows the gap at large headroom and does not change the picture (Appendix~\ref{app:v3}).

\begin{table}[t]
\caption{Regime map: Spearman correlation of $x=rd/h$ with $L$, the share of the unconstrained rule's improvement given up, and median $L$ by bin of $x$. Source: \texttt{exposure\_ratio\_test.json}.}\label{tab:regime}
\centering\scriptsize
\resizebox{\columnwidth}{!}{\begin{tabular}{@{}l r l r r r r@{}}
\toprule
 & units & Spearman$(x,L)$ [95\% CI] & \multicolumn{4}{c}{median $L$ by $x=rd/h$} \\
\cmidrule(lr){4-7}
Arm & & & $<.5$ & $.5$--$1$ & $1$--$2$ & $\ge 2$ \\
\midrule
simulation, full & 32 & 0.69 [0.55, 0.81] & 0.40 & 0.60 & 1.19 & 1.16 \\
simulation, paced & 32 & 0.73 [0.62, 0.82] & 0.11 & 0.61 & 1.17 & 1.12 \\
simulation, v3 & 32 & 0.78 [0.71, 0.83] & 0.12 & 0.81 & 1.15 & 1.14 \\
real, full & 10 & 0.83 [0.64, 0.95] & 0.33 & 0.58 & 0.90 & 1.10 \\
real, paced & 10 & 0.80 [0.63, 0.92] & 0.05 & 0.50 & 0.98 & 1.08 \\
real, v3 & 10 & 0.83 [0.67, 0.94] & 0.04 & 0.50 & 1.00 & 1.08 \\
\bottomrule
\end{tabular}
}
\end{table}

\section{Related work}
\paragraph{Conservative bandits and safe exploration.}
The closest line keeps performance above a baseline at all times. Conservative bandits require cumulative reward to stay above a fixed fraction of the baseline's, uniformly over time, in the stochastic and adversarial settings \cite{wu2016conservative}; the idea extends to contextual linear bandits \cite{kazerouni2017conservative}, reinforcement learning \cite{garcelon2020conservative}, and combinatorial semi-bandits via interleaving \cite{katariya2019conservative}. Our shield is in the same family: it deploys the risky action only if a worst-case check on the accumulated outcome passes. The differences are the following. The constraint there is multiplicative in the baseline's reward and is checked on feedback observed immediately for the played arm; here the budget is an additive amount $R$ on realised harm that includes switching fees, and the challenger's loss is learned only through purchased labels that arrive after a delay or never, so the shield pays reserves for pending feedback and for the rollback. Sampling is design-based with a label budget, which is absent there. Wu et al.\ prove regret bounds and a lower bound on the regret cost of the constraint; we give a lower bound on the reserve itself (Proposition~\ref{prop:reserves}). The most recent and closest work we found is Prudent-Banker \cite{hu2026prudent}, which treats $K$-armed adversarial bandits with arbitrary, unknown delays and bounds the expected regret against a fixed exploratory baseline by $O(\log D)$, where $D$ is the total delay. It handles unobserved feedback with a restart threshold whose slack dominates the worst-case contribution of missing rounds. Its guarantee is in expectation rather than pathwise with high probability, it has no switching fee and no rollback, feedback arrives passively for every played round rather than through a label budget with design-based sampling, and its lower bound concerns the regret rate rather than the reserve a pathwise budget requires. We found no prior work that treats this combination in our search.

\paragraph{Safe policy improvement and confidence sequences.}
High-confidence off-policy evaluation and policy improvement give a lower confidence bound on a new policy's value before it is deployed \cite{thomas2015hcope,thomas2015hcpi}. That is a one-shot certificate from logged data, with no sequential rollback and no pending feedback. Off-policy confidence sequences supply time-uniform bounds for contextual bandits and apply them to gated deployment \cite{karampatziakis2021offpolicy}; they assess the new policy's value, whereas we bound the cumulative harm of the deployed sequence, fees included. Our width uses Bennett and exact-coin mixtures in the manner of time-uniform confidence sequences \cite{howard2021time,waudbysmith2024betting}. Design-based confidence sequences for delayed outcomes monitor a counterfactual cumulative reward in randomised experiments \cite{lindon2026design}; they show that inverse-propensity error is a martingale only in a suitable event-time filtration, which our settlement accounting avoids by booking unsettled rounds at worst case. They do not control an intervention.

\paragraph{Delayed feedback and switching costs.}
Online learning under delayed feedback shows that delay multiplies regret in adversarial problems and adds to it in stochastic ones \cite{joulani2013delayed}. We do not target regret; delay enters as a reserve. Switching costs change the achievable regret rate under bandit feedback \cite{cesabianchi2013switching}, and observation budgets interact with switching costs \cite{cheng2023feedback}. These works compare to the best fixed action and impose no constraint on realised harm.

\paragraph{Label-efficient model selection and monitoring.}
Active model selection queries labels to identify a good model among pre-trained ones \cite{karimi2021online,liu2024cams}. The objective is selection or regret, not a budget on harm, and there is no switching fee, rollback or delay. Tracking the risk of a deployed model with confidence sequences flags harmful shifts and tolerates delayed labels \cite{podkopaev2022tracking}; it is an alarm, not a controller with a harm guarantee. A drift controller gates costly interventions on an online risk certificate and escalates to rollback \cite{lamaakal2026drift2act}; its certificate is on the current window risk, not on cumulative intervention harm. A challenger-switching analysis studies when to train and evaluate a challenger on a growing dataset and whether to switch, with regret against a foresight oracle \cite{digalakis2025challenger}. None gives a budget on the harm of a reversible promotion with fees under delayed, missing and sampled feedback, and none shows that the reserves are necessary.

\section{Limitations and conclusion}
The 14 real streams were used by earlier work, though no arm of this study was run or tuned on them before; each stream has few episodes, so per-stream intervals are wide or absent. The baselines are in-house heuristics or documented adapters rather than native-source reproductions. Simulated and cached outcomes are exogenous (O1), so the study does not cover actions that change outcomes. The guarantee concerns harm against never switching with a quantity $R$ the user must choose, and the label-scarce variant of the shield was designed after the real-stream run and evaluated there as a second look.

We bound the realised harm of reversible model promotion under delayed, missing and sampled feedback, with a shield that holds exactly the reserves that the matching lower bound shows are needed. The budget held in every registered and real run, the heuristics that approximate it did not, and the retention of the unconstrained rule's gain followed the exposure ratio the same bound predicts.

\bibliography{refs}
\bibliographystyle{icml2026}

\newpage
\appendix
\onecolumn
\section{Full proofs}\label{app:proofs}
\emph{Draft note (delete before submission): figures are pending. The project's figure tool (PaperBanana) was not installed when this draft was prepared, so all results are tables and one algorithm box.}

\subsection{Lemma~\ref{lem:bennett}}
Write $\xi_s=\chi_s(1-q_s/\pi_s)$ for monitored $s$. We work conditionally on the tape (losses, flags and delay mechanism, which may react to the past but not to the current coin), so the only randomness is the coins. Then $\chi_s$ and $\pi_s$ are predictable and $\E[q_s\mid\text{past}]=\pi_s$, so $\E[\xi_s\mid\text{past}]=0$ and $\mathrm{Var}(\xi_s\mid\text{past})=\chi_s^2(1/\pi_s-1)\le1/\pi_s-1$. Also $\xi_s\le b$: if $\chi_s\ge0$ then $\xi_s\le\chi_s\le1$; if $\chi_s<0$ then $\xi_s\le|\chi_s|(1/\pi_s-1)\le1/\pi_{\min}-1$. Bennett's inequality gives $\E[e^{\lambda\xi_s}\mid\text{past}]\le\exp(\phi(\lambda)(1/\pi_s-1))$, so $e^{\lambda M_t-\phi(\lambda)V_t}$ is a nonnegative supermartingale started at 1, and so is the $w$-mixture $\sum_\lambda w_\lambda e^{\lambda M_t-\phi(\lambda)V_t}$. By Ville's inequality the mixture stays below $1/\delta$ for all $t$ with probability at least $1-\delta$. On that event, for each $\lambda$, $\lambda M_t<\ln(1/(\delta w_\lambda))+\phi(\lambda)V_t$, so $M_t\le D_\delta(V_t)$. If $V_t=0$ every monitored round had $\pi=1$, hence $q=1$ and $M_t=0$. $\square$

\subsection{Pathwise bound and Theorem~\ref{thm:main}}
\emph{Bound.} For exposed rounds $s\le t$: blind rounds satisfy $\chi_s\le1$; for a monitored round $\chi_s=\chi_sq_s/\pi_s+\xi_s$, and $\chi_sq_s/\pi_s$ equals $\chi_s/\pi_s$ if queried and released, is at most $1/\pi_s$ if queried and pending, and is $0$ if not queried. Rounds that are not exposed have $\chi_s=0$. Hence $\sum_{s\le t}\chi_s\le n_{\rm b}+A+P+M_t$.

\emph{Validity argument.} Let $E$ be the event of Lemma~\ref{lem:bennett} (or Lemma~\ref{lem:coin}). Suppose at every round $t$ with $a_t=1$ the shield has verified $W_t\le R$ for a quantity with $H_t+k\le W_t$ on $E$. We claim $\sup_tH_t\le R$ on $E$. Take any $t$ and let $t'\le t$ be the last round with $a_{t'}=1$; if none exists, $H_t=0\le R$. Then $H_{t'}\le W_{t'}-k\le R-k$. If $t=t'$ we are done. Otherwise the controller rolls back at $t'+1$, paying $k$, and takes no further action change through $t$ (a later deployment would contradict the choice of $t'$), and $\chi_s=0$ for $a_s=0$. So $H_t=H_{t'}+k\le R$. No monotone invariant is needed.

\emph{The condition holds for $W_t$.} At a deployed round $t$ with proposed rate $\pi_t$, fees through round $t$ equal $k(N_{t-1}+\one\{a_{t-1}=0\})$. The bound above for rounds $\le t-1$ is $n_{\rm b}+A+P+M_{t-1}$, and round $t$ adds at most $g_t\kappa(\pi_t)$ to $n_{\rm b}+A+P$ (a blind round is at most $g_t$; a monitored round $\chi_t q_t/\pi_t\le g_t/\pi_t$) and $g_t(1/\pi_t-1)$ to $V$. On $E$, $M_t\le D_\delta(V_t)$, and $D_\delta$ is nondecreasing in $V$. Therefore $H_t\le W_t(\pi_t)-k$ on $E$. This gives $\PP(\sup_tH_t>R)\le\delta$. With $\pi\equiv1$, $V\equiv0$, $M\equiv0$ and the bound is deterministic.

\emph{Ride-out.} If $kN_{t-1}+n_{\rm b}+A+P+D_\delta(V)+(T-t+1)\le R$ while $a_{t-1}=1$, then for every $s\ge t$ on $E$ the harm is at most $kN_{t-1}$ plus the bound for the first $t-1$ rounds plus at most one unit per further round, so $H_s\le R$ with no further queries and no rollback. $\square$

\subsection{Proposition~\ref{prop:reserves}}
Let $\mathcal E'$ coincide with $\mathcal E$ on all rounds before $j$ (losses, flags, delays) and have $Z_s=-1$ on every disagreement round $s\ge j$ (so $\ell_s(0)=0$, $\ell_s(1)=1$), with the same flags. A label visible at some round $s\le j+d-1$ was requested at or before $s-d\le j-1$, where the environments coincide. With the same coins the controller therefore takes the same actions through round $j+d-1$ in both environments. Since $m\le d$, $\mathcal A$ is determined by actions through $j+m-1\le j+d-1$, and $\PP_{\mathcal E'}(\mathcal A)=\PP_{\mathcal E}(\mathcal A)>\delta$. Under $\mathcal E'$, $\chi_s\ge0$ for all $s$ and $\chi_s=1$ on disagreement rounds with $a_s=1$, so $H_t$ never decreases. On $\mathcal A$, $H_{j-1}=0$ and $H_{j+m-1}=k+e$. If the controller later rolls back at round $\rho$, $H_\rho\ge k+e+k$. If it never rolls back after $j$, $a_s=1$ for all $s\ge j$ and $H_T=k+n$. Since $n\ge e$, $\sup_tH_t\ge\min(2k+e,k+n)$ on $\mathcal A$. If $R<\min(2k+e,k+n)$ then $\PP_{\mathcal E'}(\sup_tH_t>R)\ge\PP_{\mathcal E'}(\mathcal A)>\delta$, contradicting validity. $\square$

\emph{Attainment on the all-harm tape.} For the full-sampling shield with $\pi\equiv1$, all rounds disagreement rounds with $Z=-1$, and an inner rule that always proposes $a=1$, at the first promotion round $t=n$ (the $n$-th exposed round) $W_t=2k+n$ for every delay $d$, because each earlier exposed round contributes one unit whether pending ($P$) or resolved ($A$) and the current round contributes $g_t=1$. The shield therefore deploys for exactly the rounds $n\le R-2k$, that is $\lfloor R-2k\rfloor$ of them. Proposition~\ref{prop:reserves} limits a valid controller to $e\le R-2k$ disagreement rounds in a window of length $m\le d$, so the shield attains the bound whenever $d\ge\lfloor R-2k\rfloor$. For $d<\lfloor R-2k\rfloor$ the proposition says nothing about holds longer than $d$ (labels arrive during the hold). Appendix Table~\ref{tab:prop3} reports the simulated first-spell length in $12$ cells.

\subsection{Proposition~\ref{prop:price}}
During the first spell $N=1$, so $W_t=2k+\sum_{\text{released}}\chi_s+\#\text{pending}+g_t$. With fixed delay $d$ the pending rounds and the current round are among the last $d$ rounds, so their number is at most $d$, and the released rounds are a prefix of the exposed sequence. Hence $W_t\le2k+d-S_n$ for some $n$, with $S_n=y_1+\dots+y_n$. A forced rollback needs $W_t>R$, i.e.\ $S_n<-m$ for some $n$. If $m<0$ the bound $e^{-2\mu m}\ge1$ is trivial. For $m\ge0$ and $\theta=2\mu$, Hoeffding's lemma for $y_i\in[-1,1]$ gives $\E[e^{-\theta y_i}\mid\text{past}]\le e^{-\theta\mu+\theta^2/2}=1$, so $e^{-\theta S_n}$ is a nonnegative supermartingale, and Ville's inequality gives $\PP(\inf_nS_n\le-m)\le e^{-\theta m}=e^{-2\mu m}$. For $\pm1$ steps with $\PP(y_i=+1\mid\text{past})=p'\ge p>\frac12$, $\E[(q/p)^{y_i}]=p'q/p+(1-p')p/q$ is decreasing in $p'$ and equals $1$ at $p'=p$, so $(q/p)^{S_n}$ is a supermartingale. Since $S_n$ is an integer, $S_n<-m$ means $S_n\le-(\lfloor m\rfloor+1)$, and Ville gives $((1-p)/p)^{\lfloor m\rfloor+1}$. $\square$

\subsection{Scope of the numerical checks}
The proposition-4 check uses 2000 runs per cell (per the design record), inner rule always-challenger, label budget not binding and no ride-out. In \PFourAbove{} of the \PFourCells{} cells the simulated frequency exceeds the bound at the point estimate; the largest excess is \PFourMaxZ{} binomial standard errors, so none is distinguishable from sampling noise. The forced rollbacks seen in beneficial regimes of the registered run are mostly outside the proposition, because there the label budget runs out and rounds are booked blind, and the inner rule starts further spells.

\section{Sharper bound with settlement}\label{app:v2}
\begin{lemma}[Exact coin mixture]\label{lem:coin}
For $x\in[-1,1]$ and rate $\pi$, $\E\exp(\lambda x(1-q/\pi))\le\exp(\psi(\lambda,\pi))$ with $\psi(\lambda,\pi)=\ln\max\{\pi e^{-\lambda(1/\pi-1)}+(1-\pi)e^{\lambda},\ \pi e^{\lambda(1/\pi-1)}+(1-\pi)e^{-\lambda}\}$. Moreover $\psi\ge0$ and $\psi(\lambda,1)=0$. With $\Psi_t(\lambda)=\sum_s\psi(\lambda,\pi_s)$ over the rounds included in $M$, and a finite grid with weights $w_\lambda$ (five values of $\lambda$ aimed at total variance $4,16,64,256,1024$, weights $1/5$), with probability at least $1-\delta$, for all $t$, $M_t\le D_t=\min_\lambda[\ln(1/(\delta w_\lambda))+\Psi_t(\lambda)]/\lambda$, and $D_t=0$ if no included round had $\pi<1$. Any per-round $\pi\in[\pi_{\min},1]$ is allowed.
\end{lemma}
\emph{Proof.} The function $x\mapsto\E\exp(\lambda x(1-q/\pi))$ is convex, hence maximal at $x=\pm1$, which gives the two terms. By Jensen the expectation is at least $\exp(\lambda\E[x(1-q/\pi)])=1$, so $\psi\ge0$; at $\pi=1$ the coin is deterministic, so $\psi=0$. Then $\exp(\lambda M_t-\Psi_t(\lambda))$ is a supermartingale, and the mixture and Ville's inequality finish as for Lemma~\ref{lem:bennett}. $\square$

\emph{Settlement.} Rounds with $\pi_s=1$ are accounted as before (worst case one while pending, then the exact $\chi_s$, no martingale term). For a round with $\pi_s<1$ fix a window $w\ge1$ (the declared delay if known, else a parameter). The round is \emph{unsettled} until $s+w$ and is booked at worst case $1$ regardless of its coin. At $s+w$ it settles: it enters $M$ (its $\psi$ is added to $\Psi$) and contributes $\chi_s/\pi_s$ if queried and released, $1/\pi_s$ if queried and unreleased, and $0$ if not queried. The settled set at time $t$ is the deterministic prefix $\{s\le t-w\}$ of the sampled rounds, so $M$ over settled rounds is the same martingale stopped at a deterministic time and Lemma~\ref{lem:coin} applies; unsettled rounds are bounded by $1$ pathwise. Validity does not depend on $w$. The check becomes
\begin{align*}
W_t=&\,k(N_{t-1}+\one\{a_{t-1}=0\})+k+[\text{blind}]+[\pi=1:\ \text{resolved}+\text{pending count}]\\
&+[\text{sampled, settled}:\ A_s+\text{unreleased }1/\pi]+[\text{sampled, unsettled count}]+g_t+D_t,
\end{align*}
and satisfies $H_t+k\le W_t$ on the event of Lemma~\ref{lem:coin}, so Theorem~\ref{thm:main} applies unchanged.

\emph{Label pacing.} On an exposed round the target rate is the labels remaining divided by $\max(\hat r(T-t+1),1)$, with $\hat r$ the running disagreement rate. If the target is at least $1$ the shield uses $\pi=1$; otherwise $\pi=\max(\pi_{\min},\text{target})$, or blind if no labels remain. If the paced choice fails $W_t\le R$ it escalates to $\pi=1$, and if that fails it rolls back. The blind-round variant additionally leaves a round blind when the headroom $R-W_t(\text{blind})$ is at least a margin.

\section{Protocols, seeds and tuning}\label{app:protocols}
\paragraph{Registered confirmation.}
Registered 2026-11-15 before any confirmation seed was opened; methods, arm specifications (18 frozen arms), metrics and decision rules are frozen at the registering commit, and the run was made once. Seeds 5303500000--5303500999 (1000 per regime, 30 regimes), disjoint from the earlier hypothesis blocks. Harm is $\sup_tH_t$ and a violation is $\sup_tH_t>R+10^{-9}$. Intervals are Wilson 95\% per regime ($n=1000$) and, for $\Delta J$, a seed-cluster bootstrap with 10{,}000 resamples of the equal-weight regime mean. Hypotheses: C1 zero violations for the deterministic arms; C2 per-regime Wilson upper bound at most $.05$ for the paced arms; C3 for each of five comparators the worst-regime exceedance of $R=20$ has Wilson lower bound above $.05$; C4 pooled $\Delta J$ of the $R=20$ shield against Wald+cutoff has upper bound below 3\% and against never switching below $0$; C5 pooled $\Delta J$ against Wald+cutoff over the six stress regimes has upper bound below $0$; C6 each ablation at $R=12$ has a worst-regime violation Wilson lower bound above $.05$. All six were supported. The arms were chosen after looking at two exploratory development runs (seed offsets 30000--30999 and 40000--40999 within the development block 5301000000, pilot offsets 20000--20199, tuning offsets 0--5); the six stress regimes were specified before any F3 run.

\paragraph{Real-stream protocol.}
Pre-specified and run once on 2026-11-16 (EXPOSED\_STREAMS\_FROZEN\_CONFIG: every stream had been used by earlier work). Data: cached paired 0/1 losses with 2000 events per episode: airlines (50 episodes), bank (8), criteo (50), elec2 (4), openml\_ldpa (32), rialto (16), spam (1), uci\_occupancy (4), weather (3), appliances (6), beijing\_pm25 (14), bike\_hour (6), household\_power (20), metro (16). Delays: 24 (metro, beijing\_pm25, bike\_hour), 144 (appliances, household\_power), spam 50, weather 1, elec2 48, and 100 for the remaining streams. $T=2000$, $B=200$, $c_q=.05$, $k=5$, 20 policy seeds per episode (development block offsets 60000--60019). Budgets $R=2k+h$. Violation rates use Wilson intervals over runs, which are optimistic because runs in an episode share a tape; episode-cluster bootstrap intervals (10{,}000 resamples) are used for $\Delta J$, absent for streams with fewer than six episodes.

\paragraph{Baselines and tuning.}
Comparators and the inner rule are tuned by a 32-configuration random search on six tuning seeds over the 24 original regimes, with the same budget for every method. None is a native-source reproduction: Wald+cutoff is a documented adapter of an earlier Wald/lower-confidence-bound rule with discounting and a late-label cutoff; CUSUM, the e-process gate, the CUSUM rollback and the stop-loss are in-house heuristics.

\paragraph{Exposure-ratio test.}
The test script and its prediction were committed before its first run, using existing outputs only: the label-scarce family (fresh seeds 5301071000--5301071499) and the 14 real streams (clean). A unit is a regime or stream, an arm and a headroom, kept where Wald+cutoff gains at least 10\% of $J$(never). The regime or stream is the bootstrap cluster because units of one regime share tapes.

\paragraph{Label-scarce variant (v3).}\label{app:v3}
The v2 shield with a lower pacing floor and a wider mixture grid, selected among \VThreeNCand{} candidates on a 47-regime label-scarce simulation family ($T=2000$, $B=200$, $k=5$) at seed offsets 70000--70199, frozen, then evaluated on fresh seeds 5301071000--5301071499 (scarce), 5301072000--5301072299 (original regimes) and on the real streams as a second look (the streams had been seen once with the earlier arms). The best selection score was \VThreeBest, against a higher cost for the earlier paced arm on the same family. The pacing floor barely mattered. Tables~\ref{tab:v3s}--\ref{tab:v3r}.

\section{Complete tables}
\begin{table}[h]\caption{All 18 frozen arms, pooled over 30 simulated regimes. $q_{95},q_{99}$, max are of $\sup_tH_t$.}\centering\small\begin{tabular}{@{}l r r r r r r@{}}
\toprule
Arm & $J$ & labels & $q_{95}$ & $q_{99}$ & max & deployed frac. \\
\midrule
never & 71.53 & 0.0 & 0 & 0 & 0 & 0.000 \\
wald\_cut & 68.47 & 38.8 & 23 & 37 & 59 & 0.560 \\
cusum & 68.97 & 47.0 & 29 & 34 & 117 & 0.454 \\
gate\_wald & 71.89 & 49.7 & 19 & 33 & 57 & 0.278 \\
cusum\_rb\_h4 & 68.97 & 47.8 & 16 & 31 & 68 & 0.417 \\
cusum\_rb\_h8 & 68.62 & 49.8 & 19 & 35 & 71 & 0.507 \\
stoploss\_R12 & 69.28 & 49.1 & 15 & 34 & 59 & 0.486 \\
stoploss\_R20 & 68.79 & 49.8 & 18 & 34 & 59 & 0.509 \\
stoploss\_R30 & 68.74 & 49.9 & 18 & 34 & 59 & 0.511 \\
f3\_pi1\_R12 & 70.32 & 42.1 & 12 & 12 & 12 & 0.377 \\
f3\_pi1\_R20 & 69.43 & 44.2 & 16 & 20 & 20 & 0.436 \\
f3\_pi1\_R30 & 69.25 & 44.5 & 17 & 28 & 30 & 0.455 \\
f3v2d\_R20 & 69.80 & 37.8 & 13 & 16 & 22 & 0.366 \\
f3v2d\_R30 & 68.72 & 40.5 & 16 & 21 & 31 & 0.464 \\
f3v2\_R20\_blind8 & 69.18 & 40.0 & 14 & 17 & 20 & 0.418 \\
f3\_pi1\_R12\_nopending & 70.29 & 44.1 & 14 & 29 & 46 & 0.422 \\
f3\_pi1\_R12\_norollback & 69.69 & 43.5 & 14 & 15 & 15 & 0.413 \\
f3\_pi25\_R12\_noconf & 71.74 & 28.2 & 15 & 24 & 53 & 0.323 \\
\bottomrule
\end{tabular}
\end{table}
\begin{table}[h]\caption{Exceedance rate of $R=20$ by regime (1000 runs per cell).}\centering\scriptsize\begin{tabular}{@{}l r r r r r r r@{}}
\toprule
Regime & Wald & CUSUM & Gate & CUSUM-rb & Stop & Shield pi1 & Shield paced \\
\midrule
abrupt\_beneficial & 0.001 & 0.001 & 0.001 & 0.001 & 0.004 & 0.000 & 0.000 \\
abrupt\_harmful & 0.120 & 0.041 & 0.022 & 0.020 & 0.021 & 0.000 & 0.000 \\
adverse\_informative\_delay & 0.120 & 0.051 & 0.022 & 0.011 & 0.018 & 0.000 & 0.000 \\
bait\_reversal & 0.034 & 0.003 & 0.001 & 0.000 & 0.000 & 0.000 & 0.001 \\
correlated\_errors & 0.000 & 0.000 & 0.000 & 0.000 & 0.000 & 0.000 & 0.000 \\
delayed\_onset & 0.000 & 0.003 & 0.010 & 0.003 & 0.002 & 0.000 & 0.000 \\
dep\_bait\_rho\_0.95 & 0.073 & 0.063 & 0.033 & 0.023 & 0.030 & 0.000 & 0.000 \\
dep\_harmful\_rho\_0.95 & 0.098 & 0.129 & 0.064 & 0.032 & 0.045 & 0.000 & 0.000 \\
dep\_null\_rho\_0.95 & 0.093 & 0.087 & 0.067 & 0.042 & 0.045 & 0.000 & 0.000 \\
dep\_rho\_0.0 & 0.000 & 0.001 & 0.001 & 0.000 & 0.001 & 0.000 & 0.000 \\
dep\_rho\_0.5 & 0.007 & 0.003 & 0.008 & 0.003 & 0.009 & 0.000 & 0.000 \\
dep\_rho\_0.8 & 0.031 & 0.027 & 0.023 & 0.014 & 0.024 & 0.000 & 0.000 \\
dep\_rho\_0.95 & 0.055 & 0.048 & 0.051 & 0.027 & 0.039 & 0.000 & 0.000 \\
dropout\_adverse & 0.852 & 0.684 & 0.635 & 0.830 & 0.841 & 0.000 & 0.000 \\
expensive\_switch & 0.143 & 0.987 & 0.145 & 0.160 & 0.000 & 0.000 & 0.000 \\
gradual\_beneficial & 0.000 & 0.001 & 0.002 & 0.000 & 0.003 & 0.000 & 0.000 \\
gradual\_harmful & 0.075 & 0.018 & 0.022 & 0.031 & 0.032 & 0.000 & 0.000 \\
informative\_delay & 0.001 & 0.000 & 0.001 & 0.000 & 0.001 & 0.000 & 0.000 \\
label\_dropout & 0.207 & 0.075 & 0.036 & 0.064 & 0.073 & 0.000 & 0.000 \\
label\_scarcity & 0.000 & 0.000 & 0.000 & 0.000 & 0.000 & 0.000 & 0.000 \\
late\_improvement & 0.007 & 0.007 & 0.008 & 0.009 & 0.016 & 0.000 & 0.000 \\
long\_persistent & 0.000 & 0.000 & 0.000 & 0.000 & 0.000 & 0.000 & 0.000 \\
multiple\_reversals & 0.016 & 0.081 & 0.070 & 0.020 & 0.018 & 0.000 & 0.000 \\
random\_delay & 0.001 & 0.001 & 0.000 & 0.000 & 0.002 & 0.000 & 0.000 \\
seasonal\_recurrence & 0.009 & 0.017 & 0.036 & 0.006 & 0.012 & 0.000 & 0.000 \\
short\_transient & 0.018 & 0.012 & 0.009 & 0.016 & 0.020 & 0.000 & 0.000 \\
small\_disagreement & 0.000 & 0.000 & 0.000 & 0.000 & 0.000 & 0.000 & 0.000 \\
stationary\_better & 0.000 & 0.000 & 0.000 & 0.001 & 0.002 & 0.000 & 0.000 \\
stationary\_null & 0.006 & 0.013 & 0.011 & 0.007 & 0.010 & 0.000 & 0.000 \\
stationary\_worse & 0.011 & 0.005 & 0.000 & 0.000 & 0.000 & 0.000 & 0.000 \\
\bottomrule
\end{tabular}
\end{table}
\begin{table}[h]\caption{Real streams, clean feedback, all 14 streams. ``viol./exceed'' lists runs above $R=24$ and $R=60$ for comparators and violations of the arm's own budget for budget arms.}\centering\scriptsize\begin{tabular}{@{}l r r r r r r r r@{}}
\toprule
Arm & $J$ & $\Delta J$ vs Wald & $\Delta J$ vs never & max harm & viol./exceed & labels & deployed & forced rb \\
\midrule
never & 597.89 & +194.1 & +0.0 & 0 & R24:0, R60:0 & 0.00 & 0.000 & n/a \\
wald\_cut & 403.76 & +0.0 & $-$194.1 & 175 & R24:475, R60:66 & 115.01 & 0.709 & n/a \\
cusum & 418.42 & +14.7 & $-$179.5 & 225 & R24:709, R60:132 & 141.19 & 0.665 & n/a \\
gate\_wald & 419.26 & +15.5 & $-$178.6 & 107 & R24:504, R60:132 & 149.47 & 0.599 & n/a \\
cusum\_rb\_h4 & 418.30 & +14.5 & $-$179.6 & 170 & R24:549, R60:114 & 160.61 & 0.652 & 0.84 \\
cusum\_rb\_h8 & 415.60 & +11.8 & $-$182.3 & 170 & R24:517, R60:111 & 162.75 & 0.677 & 0.27 \\
stoploss\_h14 & 419.76 & +16.0 & $-$178.1 & 84 & 311/4600 & 158.76 & 0.632 & 0.03 \\
stoploss\_h50 & 417.98 & +14.2 & $-$179.9 & 196 & 80/4600 & 162.65 & 0.674 & 0.01 \\
f3\_pi1\_h6 & 591.00 & +187.2 & $-$6.9 & 16 & 0/4600 & 102.60 & 0.174 & 1.21 \\
f3\_pi1\_h14 & 522.82 & +119.1 & $-$75.1 & 24 & 0/4600 & 133.64 & 0.417 & 1.32 \\
f3\_pi1\_h24 & 510.87 & +107.1 & $-$87.0 & 34 & 0/4600 & 142.81 & 0.477 & 1.19 \\
f3\_pi1\_h50 & 495.84 & +92.1 & $-$102.0 & 60 & 0/4600 & 151.20 & 0.530 & 0.71 \\
f3\_pi1\_h100 & 478.44 & +74.7 & $-$119.5 & 100 & 0/4600 & 153.80 & 0.565 & 0.42 \\
f3v2d\_h14 & 519.09 & +115.3 & $-$78.8 & 24 & 0/4600 & 118.77 & 0.394 & 1.09 \\
f3v2d\_h50 & 450.48 & +46.7 & $-$147.4 & 54 & 0/4600 & 141.37 & 0.600 & 0.68 \\
f3v2d\_h100 & 433.85 & +30.1 & $-$164.0 & 89 & 0/4600 & 146.21 & 0.644 & 0.28 \\
f3v2\_blind8\_h14 & 600.64 & +196.9 & +2.8 & 24 & 0/4600 & 113.41 & 0.173 & 1.63 \\
\bottomrule
\end{tabular}
\end{table}
\begin{table}[h]\caption{Real streams, adverse dropout, all 14 streams.}\centering\scriptsize\begin{tabular}{@{}l r r r r r r r r@{}}
\toprule
Arm & $J$ & $\Delta J$ vs Wald & $\Delta J$ vs never & max harm & viol./exceed & labels & deployed & forced rb \\
\midrule
never & 597.89 & +191.3 & +0.0 & 0 & R24:0, R60:0 & 0.00 & 0.000 & n/a \\
wald\_cut & 406.62 & +0.0 & $-$191.3 & 529 & R24:890, R60:447 & 121.43 & 0.831 & n/a \\
cusum & 414.03 & +7.4 & $-$183.9 & 363 & R24:889, R60:411 & 147.07 & 0.775 & n/a \\
gate\_wald & 407.44 & +0.8 & $-$190.4 & 386 & R24:644, R60:304 & 155.05 & 0.767 & n/a \\
cusum\_rb\_h4 & 416.46 & +9.8 & $-$181.4 & 529 & R24:921, R60:402 & 168.54 & 0.790 & 0.48 \\
cusum\_rb\_h8 & 415.44 & +8.8 & $-$182.4 & 529 & R24:884, R60:429 & 170.55 & 0.813 & 0.08 \\
stoploss\_h14 & 414.87 & +8.2 & $-$183.0 & 529 & 886/4600 & 169.71 & 0.810 & 0.00 \\
stoploss\_h50 & 413.50 & +6.9 & $-$184.4 & 529 & 432/4600 & 170.54 & 0.818 & 0.00 \\
f3\_pi1\_h6 & 590.98 & +184.4 & $-$6.9 & 16 & 0/4600 & 107.66 & 0.179 & 1.26 \\
f3\_pi1\_h14 & 521.74 & +115.1 & $-$76.2 & 24 & 0/4600 & 138.75 & 0.460 & 1.46 \\
f3\_pi1\_h24 & 508.98 & +102.4 & $-$88.9 & 34 & 0/4600 & 148.54 & 0.540 & 1.37 \\
f3\_pi1\_h50 & 492.82 & +86.2 & $-$105.1 & 60 & 0/4600 & 159.88 & 0.616 & 0.93 \\
f3\_pi1\_h100 & 476.27 & +69.7 & $-$121.6 & 110 & 0/4600 & 163.69 & 0.664 & 0.58 \\
f3v2d\_h14 & 518.95 & +112.3 & $-$79.0 & 24 & 0/4600 & 123.17 & 0.414 & 1.19 \\
f3v2d\_h50 & 447.24 & +40.6 & $-$150.7 & 58 & 0/4600 & 148.60 & 0.670 & 0.89 \\
f3v2d\_h100 & 431.58 & +25.0 & $-$166.3 & 101 & 0/4600 & 156.43 & 0.746 & 0.49 \\
f3v2\_blind8\_h14 & 601.12 & +194.5 & +3.2 & 24 & 0/4600 & 117.94 & 0.183 & 1.72 \\
\bottomrule
\end{tabular}
\end{table}
\begin{table}[h]\caption{Real streams, adverse delay, all 14 streams.}\centering\scriptsize\begin{tabular}{@{}l r r r r r r r r@{}}
\toprule
Arm & $J$ & $\Delta J$ vs Wald & $\Delta J$ vs never & max harm & viol./exceed & labels & deployed & forced rb \\
\midrule
never & 597.89 & +179.2 & +0.0 & 0 & R24:0, R60:0 & 0.00 & 0.000 & n/a \\
wald\_cut & 418.71 & +0.0 & $-$179.2 & 400 & R24:1017, R60:381 & 121.43 & 0.761 & n/a \\
cusum & 426.68 & +8.0 & $-$171.2 & 274 & R24:1050, R60:373 & 147.07 & 0.690 & n/a \\
gate\_wald & 421.29 & +2.6 & $-$176.6 & 321 & R24:1142, R60:439 & 155.05 & 0.728 & n/a \\
cusum\_rb\_h4 & 462.52 & +43.8 & $-$135.4 & 419 & R24:961, R60:317 & 167.56 & 0.620 & 0.65 \\
cusum\_rb\_h8 & 460.46 & +41.8 & $-$137.4 & 419 & R24:980, R60:308 & 169.29 & 0.640 & 0.43 \\
stoploss\_h14 & 447.45 & +28.8 & $-$150.4 & 419 & 979/4600 & 169.23 & 0.672 & 0.01 \\
stoploss\_h50 & 446.77 & +28.1 & $-$151.1 & 419 & 305/4600 & 169.92 & 0.677 & 0.00 \\
f3\_pi1\_h6 & 591.51 & +172.8 & $-$6.4 & 16 & 0/4600 & 107.93 & 0.179 & 1.25 \\
f3\_pi1\_h14 & 524.20 & +105.5 & $-$73.7 & 24 & 0/4600 & 138.14 & 0.447 & 1.42 \\
f3\_pi1\_h24 & 512.87 & +94.2 & $-$85.0 & 34 & 0/4600 & 148.10 & 0.529 & 1.28 \\
f3\_pi1\_h50 & 498.64 & +79.9 & $-$99.2 & 60 & 0/4600 & 156.34 & 0.578 & 0.79 \\
f3\_pi1\_h100 & 488.39 & +69.7 & $-$109.5 & 110 & 0/4600 & 160.42 & 0.603 & 0.51 \\
f3v2d\_h14 & 520.44 & +101.7 & $-$77.5 & 24 & 0/4600 & 122.92 & 0.412 & 1.22 \\
f3v2d\_h50 & 451.15 & +32.5 & $-$146.7 & 58 & 0/4600 & 146.51 & 0.653 & 0.80 \\
f3v2d\_h100 & 438.72 & +20.0 & $-$159.2 & 100 & 0/4600 & 151.33 & 0.699 & 0.40 \\
f3v2\_blind8\_h14 & 601.63 & +182.9 & +3.7 & 24 & 0/4600 & 117.94 & 0.187 & 1.75 \\
\bottomrule
\end{tabular}
\end{table}
\begin{table}[h]\caption{Real streams, clean feedback, the two harmful streams.}\centering\scriptsize\begin{tabular}{@{}l r r r r r r r r@{}}
\toprule
Arm & $J$ & $\Delta J$ vs Wald & $\Delta J$ vs never & max harm & viol./exceed & labels & deployed & forced rb \\
\midrule
never & 598.09 & $-$8.8 & +0.0 & 0 & R24:0, R60:0 & 0.00 & 0.000 & n/a \\
wald\_cut & 606.84 & +0.0 & +8.8 & 120 & R24:167, R60:9 & 122.45 & 0.198 & n/a \\
cusum & 613.23 & +6.4 & +15.2 & 225 & R24:250, R60:24 & 149.22 & 0.196 & n/a \\
gate\_wald & 605.90 & $-$0.9 & +7.8 & 91 & R24:92, R60:11 & 157.06 & 0.090 & n/a \\
cusum\_rb\_h4 & 613.80 & +7.0 & +15.7 & 170 & R24:160, R60:21 & 136.09 & 0.159 & 0.44 \\
cusum\_rb\_h8 & 612.48 & +5.6 & +14.4 & 170 & R24:169, R60:21 & 139.44 & 0.189 & 0.02 \\
stoploss\_h14 & 608.62 & +1.8 & +10.5 & 55 & 76/1160 & 136.88 & 0.139 & 0.10 \\
stoploss\_h50 & 612.47 & +5.6 & +14.4 & 170 & 21/1160 & 139.48 & 0.189 & 0.00 \\
f3\_pi1\_h6 & 608.99 & +2.2 & +10.9 & 16 & 0/1160 & 92.84 & 0.010 & 0.54 \\
f3\_pi1\_h14 & 609.09 & +2.3 & +11.0 & 24 & 0/1160 & 97.98 & 0.064 & 0.62 \\
f3\_pi1\_h24 & 607.70 & +0.9 & +9.6 & 34 & 0/1160 & 105.27 & 0.118 & 0.44 \\
f3\_pi1\_h50 & 609.42 & +2.6 & +11.3 & 58 & 0/1160 & 111.76 & 0.135 & 0.39 \\
f3\_pi1\_h100 & 613.00 & +6.2 & +14.9 & 92 & 0/1160 & 117.35 & 0.142 & 0.22 \\
f3v2d\_h14 & 609.69 & +2.9 & +11.6 & 24 & 0/1160 & 94.86 & 0.041 & 0.59 \\
f3v2d\_h50 & 608.63 & +1.8 & +10.5 & 50 & 0/1160 & 106.47 & 0.137 & 0.34 \\
f3v2d\_h100 & 610.93 & +4.1 & +12.8 & 85 & 0/1160 & 110.45 & 0.148 & 0.36 \\
f3v2\_blind8\_h14 & 610.85 & +4.0 & +12.8 & 24 & 0/1160 & 112.51 & 0.021 & 0.55 \\
\bottomrule
\end{tabular}
\end{table}
\begin{table}[h]\caption{Real streams, clean feedback, the twelve beneficial streams.}\centering\scriptsize\begin{tabular}{@{}l r r r r r r r r@{}}
\toprule
Arm & $J$ & $\Delta J$ vs Wald & $\Delta J$ vs never & max harm & viol./exceed & labels & deployed & forced rb \\
\midrule
never & 597.86 & +227.9 & +0.0 & 0 & R24:0, R60:0 & 0.00 & 0.000 & n/a \\
wald\_cut & 369.91 & +0.0 & $-$227.9 & 175 & R24:308, R60:57 & 113.77 & 0.794 & n/a \\
cusum & 385.95 & +16.0 & $-$211.9 & 118 & R24:459, R60:108 & 139.85 & 0.743 & n/a \\
gate\_wald & 388.15 & +18.2 & $-$209.7 & 107 & R24:412, R60:121 & 148.20 & 0.684 & n/a \\
cusum\_rb\_h4 & 385.71 & +15.8 & $-$212.2 & 108 & R24:389, R60:93 & 164.70 & 0.734 & 0.91 \\
cusum\_rb\_h8 & 382.79 & +12.9 & $-$215.1 & 96 & R24:348, R60:90 & 166.64 & 0.759 & 0.31 \\
stoploss\_h14 & 388.28 & +18.4 & $-$209.6 & 84 & 235/3440 & 162.41 & 0.715 & 0.01 \\
stoploss\_h50 & 385.56 & +15.7 & $-$212.3 & 196 & 59/3440 & 166.51 & 0.754 & 0.01 \\
f3\_pi1\_h6 & 588.00 & +218.1 & $-$9.9 & 16 & 0/3440 & 104.22 & 0.201 & 1.33 \\
f3\_pi1\_h14 & 508.44 & +138.5 & $-$89.4 & 24 & 0/3440 & 139.58 & 0.476 & 1.44 \\
f3\_pi1\_h24 & 494.73 & +124.8 & $-$103.1 & 34 & 0/3440 & 149.06 & 0.537 & 1.31 \\
f3\_pi1\_h50 & 476.91 & +107.0 & $-$121.0 & 60 & 0/3440 & 157.78 & 0.596 & 0.76 \\
f3\_pi1\_h100 & 456.01 & +86.1 & $-$141.8 & 100 & 0/3440 & 159.88 & 0.635 & 0.45 \\
f3v2d\_h14 & 503.99 & +134.1 & $-$93.9 & 24 & 0/3440 & 122.76 & 0.453 & 1.17 \\
f3v2d\_h50 & 424.13 & +54.2 & $-$173.7 & 54 & 0/3440 & 147.18 & 0.677 & 0.74 \\
f3v2d\_h100 & 404.34 & +34.4 & $-$193.5 & 89 & 0/3440 & 152.17 & 0.726 & 0.27 \\
f3v2\_blind8\_h14 & 598.94 & +229.0 & +1.1 & 24 & 0/3440 & 113.56 & 0.198 & 1.81 \\
\bottomrule
\end{tabular}
\end{table}
\begin{table}[h]\caption{Mean cost $J$ per stream, clean feedback, $h=50$ ($R=60$).}\centering\scriptsize\begin{tabular}{@{}l r r r r r @{}}
\toprule
Stream & never & Wald+cutoff & stop-loss & shield full & shield paced \\
\midrule
airlines & 776.80 & 781.94 & 783.11 & 781.65 & 781.70 \\
bank & 419.38 & 431.73 & 441.84 & 437.19 & 435.56 \\
criteo & 693.92 & 681.04 & 683.29 & 683.25 & 683.28 \\
elec2 & 389.50 & 394.46 & 396.08 & 395.45 & 394.52 \\
openml\_ldpa & 663.00 & 446.68 & 459.90 & 521.14 & 499.94 \\
rialto & 202.06 & 77.18 & 79.88 & 79.52 & 79.45 \\
spam & 849.00 & 212.55 & 213.85 & 472.40 & 213.30 \\
uci\_occupancy & 481.25 & 134.72 & 135.20 & 393.12 & 343.42 \\
weather & 504.00 & 438.90 & 442.18 & 442.90 & 441.52 \\
appliances & 821.83 & 611.61 & 694.99 & 836.66 & 848.32 \\
beijing\_pm25 & 664.57 & 544.22 & 599.70 & 618.52 & 605.80 \\
bike\_hour & 552.83 & 254.37 & 252.68 & 335.70 & 259.68 \\
household\_power & 632.45 & 527.58 & 550.54 & 573.18 & 558.69 \\
metro & 719.88 & 115.64 & 118.50 & 371.02 & 161.61 \\
\bottomrule
\end{tabular}
\end{table}
\begin{table}[h]\caption{OLS of $L$ on $\log x$ and the log labels per exposed round.}\label{tab:ols}\centering\small\begin{tabular}{@{}l r r@{}}
\toprule
Arm & coef. of $\log x$ & coef. of $\log$(labels per exposed round) \\
\midrule
simulation, full & 0.23 & $-$0.06 \\
simulation, paced & 0.31 & +0.10 \\
simulation, v3 & 0.28 & +0.05 \\
real, full & 0.13 & $-$0.12 \\
real, paced & 0.17 & +0.05 \\
real, v3 & 0.17 & +0.04 \\
\bottomrule
\end{tabular}
\end{table}
\begin{table}[h]\caption{Proposition~\ref{prop:reserves}: simulated first-spell length of the full-sampling shield on an all-harm tape.}\label{tab:prop3}\centering\small\begin{tabular}{@{}r r r r r l@{}}
\toprule
$d$ & $R$ & first spell & $\lfloor R-2k\rfloor$ & $\min(d,\lfloor R-2k\rfloor)$ & attained \\
\midrule
1 & 8 & 2 & 2 & 1 & yes \\
1 & 12 & 6 & 6 & 1 & yes \\
1 & 20 & 14 & 14 & 1 & yes \\
3 & 8 & 2 & 2 & 2 & yes \\
3 & 12 & 6 & 6 & 3 & yes \\
3 & 20 & 14 & 14 & 3 & yes \\
5 & 8 & 2 & 2 & 2 & yes \\
5 & 12 & 6 & 6 & 5 & yes \\
5 & 20 & 14 & 14 & 5 & yes \\
10 & 8 & 2 & 2 & 2 & yes \\
10 & 12 & 6 & 6 & 6 & yes \\
10 & 20 & 14 & 14 & 10 & yes \\
\bottomrule
\end{tabular}
\end{table}
\begin{table}[h]\caption{Proposition~\ref{prop:price}: simulated forced-rollback frequency against the ruin bound (2000 runs per cell); $z$ is the excess over the bound in binomial standard errors, zero when below.}\label{tab:prop4}\centering\small\begin{tabular}{@{}r r r r r r r@{}}
\toprule
$p$ & $d$ & $R$ & $m$ & forced rollbacks & bound & $z$ \\
\midrule
0.6 & 1 & 10 & +3 & 0.2035 & 0.1975 & 0.7 \\
0.6 & 1 & 12 & +5 & 0.0860 & 0.0878 & 0.0 \\
0.6 & 1 & 15 & +8 & 0.0220 & 0.0260 & 0.0 \\
0.6 & 1 & 20 & +13 & 0.0010 & 0.0034 & 0.0 \\
0.6 & 5 & 10 & $-$1 & 0.3970 & 1.0000 & 0.0 \\
0.6 & 5 & 12 & +1 & 0.1875 & 0.4444 & 0.0 \\
0.6 & 5 & 15 & +4 & 0.0530 & 0.1317 & 0.0 \\
0.6 & 5 & 20 & +9 & 0.0045 & 0.0173 & 0.0 \\
0.7 & 1 & 10 & +3 & 0.0345 & 0.0337 & 0.2 \\
0.7 & 1 & 12 & +5 & 0.0060 & 0.0062 & 0.0 \\
0.7 & 1 & 15 & +8 & 0.0005 & 0.0005 & 0.0 \\
0.7 & 1 & 20 & +13 & 0.0000 & 0.0000 & 0.0 \\
0.7 & 5 & 10 & $-$1 & 0.1670 & 1.0000 & 0.0 \\
0.7 & 5 & 12 & +1 & 0.0310 & 0.1837 & 0.0 \\
0.7 & 5 & 15 & +4 & 0.0025 & 0.0145 & 0.0 \\
0.7 & 5 & 20 & +9 & 0.0000 & 0.0002 & 0.0 \\
0.8 & 1 & 10 & +3 & 0.0045 & 0.0039 & 0.4 \\
0.8 & 1 & 12 & +5 & 0.0005 & 0.0002 & 0.7 \\
0.8 & 1 & 15 & +8 & 0.0000 & 0.0000 & 0.0 \\
0.8 & 1 & 20 & +13 & 0.0000 & 0.0000 & 0.0 \\
0.8 & 5 & 10 & $-$1 & 0.0685 & 1.0000 & 0.0 \\
0.8 & 5 & 12 & +1 & 0.0060 & 0.0625 & 0.0 \\
0.8 & 5 & 15 & +4 & 0.0005 & 0.0010 & 0.0 \\
0.8 & 5 & 20 & +9 & 0.0000 & 0.0000 & 0.0 \\
\bottomrule
\end{tabular}
\end{table}
\begin{table}[h]\caption{Label-scarce variant v3, scarce family (47 regimes), fresh seeds.}\label{tab:v3s}\centering\scriptsize\begin{tabular}{@{}l r l l l@{}}
\toprule
Arm & $J$ & $\Delta J$ vs Wald & $\Delta J$ vs never & violations \\
\midrule
never & 433.81 & +92.98 [+92.64,+93.28] & +0.00 [+0.00,+0.00] & R24:0.000; R60:0.000; R110:0.000 \\
wald\_cut & 340.83 & +0.00 [+0.00,+0.00] & -92.98 [-93.30,-92.67] & R24:0.067; R60:0.004; R110:0.000 \\
stoploss\_h50 & 364.01 & +23.18 [+22.74,+23.63] & -69.79 [-70.29,-69.33] & 1374/23500 \\
f3\_pi1\_h50 & 392.90 & +52.07 [+51.66,+52.49] & -40.90 [-41.27,-40.53] & 0/23500 \\
f3v2d\_h50 & 387.65 & +46.82 [+46.29,+47.34] & -46.15 [-46.66,-45.66] & 0/23500 \\
f3\_pi1\_h100 & 382.59 & +41.76 [+41.37,+42.16] & -51.21 [-51.58,-50.84] & 0/23500 \\
f3v2d\_h100 & 357.55 & +16.72 [+16.33,+17.10] & -76.26 [-76.66,-75.85] & 0/23500 \\
v3\_h50 & 394.86 & +54.03 [+53.39,+54.67] & -38.95 [-39.59,-38.30] & 0/23500 \\
v3\_h100 & 355.31 & +14.48 [+13.99,+14.99] & -78.49 [-79.00,-77.98] & 0/23500 \\
\bottomrule
\end{tabular}
\end{table}
\begin{table}[h]\caption{Label-scarce variant v3, original 30 regimes, fresh seeds.}\label{tab:v3o}\centering\scriptsize\begin{tabular}{@{}l r l l l@{}}
\toprule
Arm & $J$ & $\Delta J$ vs Wald & $\Delta J$ vs never & violations \\
\midrule
wald\_cut & 68.33 & +0.00 [+0.00,+0.00] & -3.25 [-3.46,-3.03] & R12:0.180; R20:0.065; R30:0.027 \\
f3\_pi1\_R12 & 70.35 & +2.02 [+1.81,+2.24] & -1.22 [-1.44,-1.01] & 0/9000 \\
f3\_pi1\_R20 & 69.50 & +1.16 [+0.97,+1.35] & -2.08 [-2.31,-1.86] & 0/9000 \\
f3\_pi1\_R30 & 69.32 & +0.99 [+0.80,+1.17] & -2.26 [-2.49,-2.04] & 0/9000 \\
f3v2d\_R20 & 69.95 & +1.62 [+1.41,+1.83] & -1.63 [-1.85,-1.41] & 0/9000 \\
f3v2d\_R30 & 68.80 & +0.47 [+0.28,+0.67] & -2.77 [-3.00,-2.55] & 0/9000 \\
v3\_R12 & 73.07 & +4.74 [+4.49,+4.98] & +1.49 [+1.28,+1.70] & 0/9000 \\
v3\_R20 & 69.95 & +1.62 [+1.41,+1.83] & -1.63 [-1.84,-1.41] & 0/9000 \\
v3\_R30 & 68.82 & +0.49 [+0.30,+0.69] & -2.75 [-2.98,-2.53] & 0/9000 \\
\bottomrule
\end{tabular}
\end{table}
\begin{table}[h]\caption{Label-scarce variant v3 on the real streams (second look), clean feedback, pooled.}\label{tab:v3r}\centering\scriptsize\begin{tabular}{@{}l r r r r@{}}
\toprule
Arm & $J$ & $\Delta J$ vs Wald & $\Delta J$ vs never & viol. \\
\midrule
f3\_pi1\_h50 & 495.84 & +92.1 & $-$102.0 & 0/4600 \\
f3v2d\_h50 & 450.48 & +46.7 & $-$147.4 & 0/4600 \\
v3\_h50 & 444.59 & +40.8 & $-$153.3 & 0/4600 \\
f3\_pi1\_h100 & 478.44 & +74.7 & $-$119.5 & 0/4600 \\
f3v2d\_h100 & 433.85 & +30.1 & $-$164.0 & 0/4600 \\
v3\_h100 & 421.40 & +17.6 & $-$176.5 & 0/4600 \\
\bottomrule
\end{tabular}
\end{table}

\end{document}
